\section{Materials and methods}

\subsection{Case studies}

We used nine semi-synthetic scenarios: three study areas (North Island, Kamloops, and Prince
George timber supply areas, BC) crossed with three problem sizes (6, 18, and 40 harvest blocks;
Table~\ref{tab:scenarios}) \citep{jaffray2026zenodo}. Block geometry, area, and stand attributes
come from an overlay of the provincial Forest Harvest Authorizations and Vegetation Resources
Inventory layers \citep{bcgov_fha,bcgov_vri}. Average ground slope per block was derived from a
digital elevation model. Every block is harvested with the same ground-based system: feller-buncher,
grapple skidder, processor, then loader. Machine productivity (m$^3$\,day$^{-1}$, per machine and
block) was predicted from stand and terrain attributes with the productivity models bundled with
FHOPS \pending{cite Lahrsen 2025, ADV6N7, Berry 2019, and Barko 450 TN-46 sources}, with a floor of
1~m$^3$\,day$^{-1}$. Fleets of 6, 12, and 24 machines were split as evenly as possible across the
four roles. Each block has its own landing, which at most two machines can work at once. The
planning horizon is 112 days with three 8-h shifts per day, 7 days per week, and full machine
availability. Blocks can be harvested at any time within the horizon. The scenarios differ in
total volume relative to fleet capacity, so in some scenarios not all volume can be harvested
within 112 days (Table~\ref{tab:scenarios}). Delivered volume, rather than completion time, is
therefore the primary measure of plan quality.

\begin{table}[htbp]
  \centering
  \caption{Scenario characteristics. Machines are given as total
  (feller-buncher/grapple skidder/processor/loader). Volume is the total volume to be harvested;
  slope and stem size are block means. All scenarios have one landing per block (capacity two
  machines), a 112-day horizon, and three 8-h shifts per day.}
  \label{tab:scenarios}
  \small
  \setlength{\tabcolsep}{4pt}
  % Generated by scripts/make_table_scenarios.py; do not edit.
\begin{tabular}{lrlrrrr}
\toprule
Study area & Blocks & Machines & Volume (m$^3$) & Mean block (m$^3$) & Slope (\%) & Stem (m$^3$) \\
\midrule
North Island & 6 & 6 (2/2/1/1) & 323,091 & 53,848 & 43.6 & 0.66 \\
North Island & 18 & 12 (3/3/3/3) & 755,750 & 41,986 & 39.9 & 0.91 \\
North Island & 40 & 24 (6/6/6/6) & 1,578,930 & 39,473 & 45.0 & 0.93 \\
Kamloops & 6 & 6 (2/2/1/1) & 30,913 & 5,152 & 21.2 & 0.23 \\
Kamloops & 18 & 12 (3/3/3/3) & 71,386 & 3,966 & 15.8 & 0.24 \\
Kamloops & 40 & 24 (6/6/6/6) & 393,614 & 9,840 & 21.0 & 0.32 \\
Prince George & 6 & 6 (2/2/1/1) & 67,924 & 11,321 & 9.0 & 0.46 \\
Prince George & 18 & 12 (3/3/3/3) & 557,762 & 30,987 & 14.5 & 0.54 \\
Prince George & 40 & 24 (6/6/6/6) & 1,309,443 & 32,736 & 14.5 & 0.56 \\
\bottomrule
\end{tabular}

\end{table}

\subsection{Scheduling model and solution methods}

FHOPS represents each scenario as a shift-level machine scheduling problem: which machine works on
which block in each shift. The objective maximizes delivered volume, subject to machine
availability, landing capacity, and the precedence between harvesting phases. A downstream phase
can only process volume that the upstream phase has already staged at the landing. The
mixed-integer linear programming (MIP) formulation is documented with the software
\citep{fhops2026formulation} and in \fhopspaper; we do not restate it here. The scenarios carry no
mobilisation distances, so moving a machine between blocks has no cost. \pending{Confirm the
objective terms active in these scenarios (production weight 1.0; mobilisation weight 0.5 has no
effect without distances).}

Each window was solved with one of two methods. (i) Simulated annealing (SA), the FHOPS
metaheuristic: 500 iterations per window, random seed 42. (ii) The FHOPS operational MIP, solved
with \mipbackend\ from a cold start (no heuristic warm start) on one thread, with a limit of
\miptimelimit\ per window. In rolling windows FHOPS also runs a lexicographic second stage. It
favours earlier production without worsening the first-stage objective, so that a window does not
postpone work beyond its locked span \citep{fhops2026v101}. This stage was limited to
\earlinesslimit. \pending{Justify the limits with the time-limit calibration
(incumbent/bound trajectories), and state whether a deterministic work limit was used.}

\subsection{Rolling-horizon procedure with state carry-forward}

For a master horizon of $H = 112$ days, a planning window of $\theta$ days, and a re-optimization
interval (lock span) of $f \le \theta$ days, FHOPS~1.0.1 proceeds as follows
\citep{fhops2026v101}:
\begin{enumerate}
  \item Solve the subproblem for days $[t, \min(t+\theta-1, H)]$, starting at $t=1$.
  \item Lock the decisions for days $[t, t+f-1]$ and append them to the stitched plan.
  \item Replay the stitched plan with the deterministic playback engine to obtain the carried
        state: the volume remaining in each block, the remaining output, staged inventory, and
        shift counts of each phase, and the last block of each machine.
  \item Set $t \leftarrow t+f$ and repeat from that state until $t > H$.
\end{enumerate}
A window for which the solver returns no solution locks nothing for its span and is recorded as
such. A window with nothing left to plan is recorded as skipped. A window that returns a plan
delivering no volume is recorded as empty. Earlier FHOPS releases reset every window to the
initial scenario state, and \citet{jaffray2026thesis} used one of them. Re-planning the full block
volume in every window, with staged inventory reset, inflated the apparent effect of the
re-optimization interval. The results reported here supersede those values.

\subsection{Experimental design}

The full factorial design crossed the nine scenarios with
$\theta \in \{2, 4, 8, 16\}$~weeks, $f \in \{1, 7, 14\}$~days, and both solution methods: 216
rolling-horizon runs. These levels represent daily, weekly, and biweekly re-planning cycles and
horizons from two weeks to the full season. For each scenario and method we also solved one
full-horizon baseline ($\theta = f = H$), giving 18 baselines. Inputs were identical across
windows, so differences between runs come only from the rolling-horizon configuration. Runs were
executed in parallel, one process per run, \pending{hardware and concurrency}.

\subsection{Performance measures}

Each run was scored on its stitched plan, replayed against the full 112-day scenario with FHOPS
deterministic playback. We report:
\begin{itemize}
  \item delivered volume (m$^3$ output by the final phase) and remaining volume;
  \item relative delivered volume: the run's delivered volume divided by that of the reference
        plan, either the same-method full-horizon baseline or the best stitched plan found for
        the scenario \pending{choose reference};
  \item sequencing violations in the replayed plan, counted in total and for those whose volume
        exceeds \pending{threshold}~m$^3$;
  \item block-to-block machine moves and schedule continuity (the share of block--machine
        assignments unchanged between consecutive working days of a block), computed from the
        stitched assignments;
  \item computing effort: total wall time over all windows, and the number of windows that
        reached the limit, returned no solution, or were empty.
\end{itemize}
The design is deterministic, with one run per configuration, so we compare configurations
descriptively rather than with inferential tests.
